\documentclass[11pt]{article}
\usepackage[a4paper,margin=0.85in]{geometry}
\usepackage{amsmath,amssymb,mathtools,bm}
\usepackage{booktabs}
\usepackage{array}
\usepackage{tabularx}
\usepackage{xcolor}
\usepackage[colorlinks=true,linkcolor=blue!45!black,urlcolor=blue!45!black]{hyperref}

\setlength{\parindent}{0pt}
\setlength{\parskip}{0.55em}
\setlength{\emergencystretch}{3em}
\renewcommand{\arraystretch}{1.15}
\newcolumntype{Y}{>{\raggedright\arraybackslash}X}

\newcommand{\E}{\mathbb E}
\newcommand{\TC}{T^C}
\newcommand{\TN}{T^N}
\newcommand{\TY}{T^Y}
\newcommand{\tauc}{\tau^C}
\newcommand{\taun}{\tau^N}
\newcommand{\tauy}{\tau^Y}
\newcommand{\bg}{b^g}
\newcommand{\ig}{i^g}
\newcommand{\Pc}{P^C}
\newcommand{\MC}{MC}
\newcommand{\Rev}{\operatorname{Rev}}

\title{Fiscal-LPT Appendix Derivation Audit}
\author{Adversarial derivation check}
\date{May 13, 2026}

\begin{document}
\maketitle

\section{Scope and verdict}

This note audits the appendix derivations in \path{Fiscal-LPT/main.tex} against the
live Julia implementation in \path{MCMS-private/src/build_equations.jl} and the
current calibration block in \path{MCMS-private/src/calibration.jl}. The equations
implemented in Julia are internally coherent under a precise set of conventions.
The appendix mostly reports the implemented equations correctly, but several
derivation statements are too loose. One item is a real derivation fault: the
level government budget constraint written in the appendix is not the correct
level ancestor of the implemented linear budget equation.

\section{Adversarial findings}

\textbf{Fix: opening appendix conventions.}
The text says the new fiscal objects are log deviations. This is not true for
\(\bg_{k,t}\). The implemented debt variable is an additive real public-debt
deviation normalized by steady-state GDP, not a log deviation. Treating debt as
a log deviation would change the fiscal-rule and budget-constraint scaling.

\textbf{Fix: labour-tax derivation and revenue equation.}
The appendix must define \(\taun_{k,t}\equiv \widehat{(1-\TN_{k,t})^{-1}}\), not
just ``the log labour wedge.'' The wage Phillips curve and the labour-tax
revenue equation both require this inverse after-tax wedge convention. If
\(\taun\) were instead \(\widehat{1-\TN}\), the signs would flip.

\textbf{Fix: government budget derivation.}
The appendix starts from
\[
  B^g_{t}=(1+i^g_{t-1})B^g_{t-1}-P^C_tY_t\Rev_t.
\]
That timing and gross-debt form do not linearize to the implemented equation
\[
  \beta \bg_t-\beta\bar b^g\ig_t+\Rev_t=\bg_{t-1}-\bar b^g\pi^C_t.
\]
The implemented equation comes from a one-period discount-debt budget,
\[
  Q^g_tB^g_t + P^C_t\bar Y\,\Rev_t = B^g_{t-1},
  \qquad Q^g_t=(R^g_t)^{-1}.
\]
This should be corrected in the appendix.

\textbf{Clarify: revenue accounting.}
The revenue variables \(rev^C,rev^N,rev^Y,rev^{tar}\) are not log revenues.
They are first-order deviations of revenue scaled by steady-state GDP. The
formulas are correct only under that additive normalization.

\textbf{Clarify: consumption-tax Euler derivation.}
The Euler equation is correct, but the text should say that \(i_{k,t}\) is the
log deviation of the gross nominal rate, or the model's first-order interest-rate
gap measured in the same units as inflation. Otherwise the coefficient on the
interest-rate term is ambiguous.

\textbf{Clarify: production-tax and subsidy experiments.}
A statement such as \(\sigma^Y_{4,21}=-0.10\) is a shock-scale override. It is a
ten percentage point subsidy only for a unit innovation and under the local
mapping \(d\TY=(1+\bar \TY)\tauy\). If the model shock is interpreted as a
statutory-rate movement, the gross-wedge conversion must be stated.

\textbf{Verify: Dynare mirror claim.}
\texttt{main.tex} says there is a Dynare mirror in \texttt{b*\_fiscal\_lpt.mod}.
The model equations in the mirror match the structure, but the parameter file
still contains older placeholder assignments such as \(\bar \TC=0.20\),
\(\bar \TN=0.30\), \(\bar b^g=0.60\), and \(\psi^g=0.01\). The current Julia
calibration instead uses country-specific values and \(\psi^g=0\). Do not cite
the Dynare parameter file as the current quantitative calibration unless it is
regenerated or patched.

\section{Notation that makes the derivations valid}

For strictly positive variables, lower-case variables are log deviations:
\[
  x_t=\widehat X_t=\log X_t-\log \bar X.
\]
Debt is different. The implemented debt state is an additive deviation of real
public debt from its steady-state ratio:
\[
  \bg_{k,t}
  =
  \frac{B^g_{k,t}}{\Pc_{k,t}\bar Y_k}-\bar b^g_k.
\]
It is normalized by current CPI and steady-state real GDP. This convention is
forced by the implemented budget equation because no current \(y_{k,t}\) term
appears there.

The three tax variables are log deviations of gross wedges:
\begin{align}
  \tauc_{k,t} &\equiv \widehat{(1+\TC_{k,t})}, \\
  \taun_{k,t} &\equiv \widehat{(1-\TN_{k,t})^{-1}}, \\
  \tauy_{k i,t} &\equiv \widehat{(1+\TY_{k i,t})}.
\end{align}
Therefore the first-order statutory-rate changes are
\begin{align}
  d\TC_{k,t} &= (1+\bar\TC_k)\tauc_{k,t}, \\
  d\TN_{k,t} &= (1-\bar\TN_k)\taun_{k,t}, \\
  d\TY_{k i,t} &= (1+\bar\TY_{k i})\tauy_{k i,t}.
\end{align}
These conversion formulas are the hidden step behind the revenue equations.

\section{Consumption tax and the Euler equation}

The household consumption FOC is
\[
  U_C(C_{k,t})=\Lambda_{k,t}\Pc_{k,t}(1+\TC_{k,t}),
\]
so, with CRRA utility \(U_C(C)=C^{-\sigma}\),
\[
  \log \Lambda_{k,t}
  =
  -\sigma \log C_{k,t}
  -\log \Pc_{k,t}
  -\log(1+\TC_{k,t}).
\]
After subtracting the steady state,
\[
  \lambda_{k,t}=-\sigma c_{k,t}-p^C_{k,t}-\tauc_{k,t}.
\]

The nominal bond FOC is
\[
  \Lambda_{k,t}
  =
  \beta R_{k,t}\E_t\Lambda_{k,t+1},
  \qquad R_{k,t}=1+i^{level}_{k,t}.
\]
At first order, using \(\beta \bar R=1\) and ignoring Jensen terms,
\[
  \lambda_{k,t}
  =
  i_{k,t}+\E_t\lambda_{k,t+1},
\]
where \(i_{k,t}\) is the log deviation of the gross nominal rate. Substitute the
multiplier expression:
\[
  -\sigma c_{k,t}-p^C_{k,t}-\tauc_{k,t}
  =
  i_{k,t}
  +\E_t[-\sigma c_{k,t+1}-p^C_{k,t+1}-\tauc_{k,t+1}].
\]
Move terms and use \(\pi^C_{k,t+1}=p^C_{k,t+1}-p^C_{k,t}\):
\[
  \sigma c_{k,t}
  =
  \sigma \E_t c_{k,t+1}
  -i_{k,t}
  +\E_t\pi^C_{k,t+1}
  +\E_t\tauc_{k,t+1}
  -\tauc_{k,t}.
\]
Divide by \(\sigma\):
\[
  c_{k,t}
  =
  \E_t c_{k,t+1}
  -\frac{1}{\sigma}
  \left(
    i_{k,t}
    -\E_t\pi^C_{k,t+1}
    -\E_t\tauc_{k,t+1}
    +\tauc_{k,t}
  \right).
\]
This exactly matches the implemented Julia equation
\[
  c_k[0]=c_k[1]-\sigma^{-1}(i_k[0]-\pi^C_k[1]-\tauc_k[1]+\tauc_k[0]).
\]

\section{Labour tax and the wage Phillips curve}

The static labour condition is
\[
  \frac{W_{k,t}}{\Pc_{k,t}}
  =
  \frac{1+\TC_{k,t}}{1-\TN_{k,t}}
  C_{k,t}^{\sigma}N_{k,t}^{\varphi}.
\]
Taking deviations gives
\[
  w^{flex}_{k,t}
  =
  \sigma c_{k,t}
  +\varphi n_{k,t}
  +\tauc_{k,t}
  +\taun_{k,t}.
\]
The sign on \(\taun\) is positive because \(\taun\) is the log deviation of the
inverse after-tax labour wedge. A higher labour tax raises the pre-tax wage
required by households.

The Calvo wage block uses the desired real-wage gap:
\[
  \pi^w_{k,t}
  =
  \kappa^w_k(w^{flex}_{k,t}-w_{k,t})
  +\beta\E_t\pi^w_{k,t+1}.
\]
Substitution yields
\[
  \pi^w_{k,t}
  =
  \kappa^w_k
  \left(
    \sigma c_{k,t}
    +\varphi n_{k,t}
    -w_{k,t}
    +\tauc_{k,t}
    +\taun_{k,t}
  \right)
  +\beta\E_t\pi^w_{k,t+1}.
\]
This matches the implemented wage equation. The only appendix problem is
definition: without the inverse after-tax wedge convention, the sign is not
well-defined.

\section{Production tax and price setting}

The selected instrument is a marginal-cost wedge:
\[
  \MC^{tax}_{k i,t}=(1+\TY_{k i,t})\MC_{k i,t}.
\]
Taking log deviations,
\[
  mc^{tax}_{k i,t}
  =
  mc_{k i,t}+\tauy_{k i,t}.
\]
Hence every price Phillips-curve marginal-cost gap receives the same
\(+\tauy_{k i,t}\) term. For producer-currency pricing:
\[
  \pi_{lki,t}
  =
  \kappa^p_{ki}(mc_{ki,t}+\tauy_{ki,t}-p_{lki,t})
  +\beta\E_t\pi_{lki,t+1}.
\]
For local-currency pricing with \(l\ne k\), the implemented relative-price term
also includes the bilateral real exchange rate:
\[
  \pi_{lki,t}
  =
  \kappa^p_{ki}(mc_{ki,t}+\tauy_{ki,t}-p_{lki,t}-q_{kl,t})
  +\beta\E_t\pi_{lki,t+1}.
\]
For dominant-currency pricing, the \(q\) term is relative to the dominant
currency country. The appendix is acceptable because it says it suppresses the
currency-pricing indicators, but the full derivation should mention these
branches if the appendix is meant to be self-contained.

\section{Fiscal rules}

For a generic wedge \(x_{k,t}\), define a target
\[
  x^{target}_{k,t}
  =
  \phi^x_{y,k}y_{k,t}
  +
  \phi^x_{b,k}\bg_{k,t-1}.
\]
Partial adjustment gives
\[
  x_{k,t}
  =
  \rho^x_k x_{k,t-1}
  +(1-\rho^x_k)x^{target}_{k,t}
  +\sigma^x_k\varepsilon^x_{k,t}.
\]
Therefore
\begin{align}
  \tauc_{k,t}
  &=
  \rho^C_k\tauc_{k,t-1}
  +(1-\rho^C_k)(\phi^C_{y,k}y_{k,t}+\phi^C_{b,k}\bg_{k,t-1})
  +\sigma^C_k\varepsilon^C_{k,t}, \\
  \taun_{k,t}
  &=
  \rho^N_k\taun_{k,t-1}
  +(1-\rho^N_k)(\phi^N_{y,k}y_{k,t}+\phi^N_{b,k}\bg_{k,t-1})
  +\sigma^N_k\varepsilon^N_{k,t}, \\
  \tauy_{k i,t}
  &=
  \rho^Y_{k i}\tauy_{k i,t-1}
  +(1-\rho^Y_{k i})(\phi^Y_{y,k i}y_{k,t}+\phi^Y_{b,k i}\bg_{k,t-1})
  +\sigma^Y_{k i}\varepsilon^Y_{k i,t}.
\end{align}
These equations match the Julia implementation. The only convention issue is
again that \(\bg\) is an additive normalized debt deviation, not a log deviation.

\section{Revenue linearizations}

The implemented revenue variables are additive deviations normalized by
steady-state GDP. They are not logs.

\subsection{Consumption-tax revenue}

Let
\[
  s^C_k=\frac{\bar\Pc_k\bar C_k}{\bar\Pc_k\bar Y_k}
  =\frac{\bar C_k}{\bar Y_k}.
\]
Revenue is
\[
  \mathcal R^C_{k,t}=\TC_{k,t}\Pc_{k,t}C_{k,t}.
\]
Scaled by steady-state nominal GDP and abstracting from the common CPI level
already handled in the real budget normalization,
\[
  rev^C_{k,t}
  =
  s^C_k
  \left[
    d\TC_{k,t}
    +\bar\TC_k c_{k,t}
  \right].
\]
Using \(d\TC=(1+\bar\TC)\tauc\),
\[
  rev^C_{k,t}
  =
  s^C_k
  \left[
    (1+\bar\TC_k)\tauc_{k,t}
    +\bar\TC_k c_{k,t}
  \right].
\]

\subsection{Labour-tax revenue}

Let
\[
  s^L_k=\frac{\bar W_k\bar N_k}{\bar\Pc_k\bar Y_k}.
\]
Revenue is
\[
  \mathcal R^N_{k,t}=\TN_{k,t}W_{k,t}N_{k,t}.
\]
The first-order scaled deviation is
\[
  rev^N_{k,t}
  =
  s^L_k
  \left[
    d\TN_{k,t}
    +\bar\TN_k(w_{k,t}+n_{k,t})
  \right].
\]
Using \(d\TN=(1-\bar\TN)\taun\),
\[
  rev^N_{k,t}
  =
  s^L_k
  \left[
    (1-\bar\TN_k)\taun_{k,t}
    +\bar\TN_k(w_{k,t}+n_{k,t})
  \right].
\]

\subsection{Production-tax revenue}

Let
\[
  s^{MCY}_{k i}=\frac{\overline{\MC}_{k i}\bar Y_{k i}}{\bar\Pc_k\bar Y_k}.
\]
Revenue is
\[
  \mathcal R^Y_{k,t}=\sum_i \TY_{k i,t}\MC_{k i,t}Y_{k i,t}.
\]
The first-order scaled deviation is
\[
  rev^Y_{k,t}
  =
  \sum_i s^{MCY}_{k i}
  \left[
    d\TY_{k i,t}
    +\bar\TY_{k i}(mc_{k i,t}+y_{k i,t})
  \right].
\]
Using \(d\TY=(1+\bar\TY)\tauy\),
\[
  rev^Y_{k,t}
  =
  \sum_i s^{MCY}_{k i}
  \left[
    (1+\bar\TY_{k i})\tauy_{k i,t}
    +\bar\TY_{k i}(mc_{k i,t}+y_{k i,t})
  \right].
\]
This matches the implemented equation. If \(\bar\TY<0\), the base-effect term
has the sign of a subsidy: higher marginal-cost output makes the subsidy more
expensive.

\section{Government borrowing spread}

The spread equation is directly linear:
\[
  \ig_{k,t}=i_{k,t}+\psi^g_k\bg_{k,t-1}.
\]
The baseline calibration currently sets \(\psi^g_k=0\) in Julia. Therefore the
spread channel is present as an equation but shut down quantitatively in the
current Julia baseline.

\section{Government budget constraint}

The implemented budget is a discount-debt equation. Start from
\[
  Q^g_{k,t}B^g_{k,t}
  +
  \Pc_{k,t}\bar Y_k
  \left(
    rev^C_{k,t}+rev^N_{k,t}+rev^Y_{k,t}+rev^{tar}_{k,t}
  \right)
  =
  B^g_{k,t-1}.
\]
Divide by \(\Pc_{k,t}\bar Y_k\):
\[
  Q^g_{k,t}
  \frac{B^g_{k,t}}{\Pc_{k,t}\bar Y_k}
  +rev^C_{k,t}+rev^N_{k,t}+rev^Y_{k,t}+rev^{tar}_{k,t}
  =
  \frac{B^g_{k,t-1}}{\Pc_{k,t}\bar Y_k}.
\]
Use
\[
  \frac{B^g_{k,t}}{\Pc_{k,t}\bar Y_k}=\bar b^g_k+\bg_{k,t},
  \qquad
  \frac{B^g_{k,t-1}}{\Pc_{k,t}\bar Y_k}
  =
  \frac{\Pc_{k,t-1}}{\Pc_{k,t}}
  \frac{B^g_{k,t-1}}{\Pc_{k,t-1}\bar Y_k}.
\]
At first order,
\[
  \frac{\Pc_{k,t-1}}{\Pc_{k,t}}(\bar b^g_k+\bg_{k,t-1})
  =
  \bar b^g_k+\bg_{k,t-1}-\bar b^g_k\pi^C_{k,t}.
\]
For discount debt, \(Q^g_{k,t}=1/R^g_{k,t}\). With \(\bar Q^g=\beta\) and
\(\ig_{k,t}\) as the log deviation of the gross government rate,
\[
  Q^g_{k,t}(\bar b^g_k+\bg_{k,t})
  =
  \beta\bar b^g_k+\beta\bg_{k,t}-\beta\bar b^g_k\ig_{k,t}.
\]
Subtract the steady-state identity \(\beta\bar b^g_k=\bar b^g_k-\bar rev_k\).
The first-order budget is
\[
  \beta\bg_{k,t}
  -\beta\bar b^g_k\ig_{k,t}
  +rev^C_{k,t}+rev^N_{k,t}+rev^Y_{k,t}+rev^{tar}_{k,t}
  =
  \bg_{k,t-1}
  -\bar b^g_k\pi^C_{k,t}.
\]
This is the implemented equation. The current appendix should use this
discount-debt parent equation. Starting from a lagged gross-interest budget is
not just a harmless alternative presentation; it changes the timing and the
linear coefficients.

\section{Minimum edits recommended for \texttt{main.tex}}

The appendix can be made mathematically tight with these targeted edits:
\begin{enumerate}
  \item Define \(\tauc,\taun,\tauy\) exactly as gross-wedge log deviations and
  state that \(\bg\) is an additive debt-ratio deviation.
  \item Replace the government budget parent equation with
  \(Q^g_tB^g_t+\Pc_t\bar Y\,rev_t=B^g_{t-1}\).
  \item Say explicitly that revenue variables are first-order GDP-normalized
  deviations, not log revenue variables.
  \item Qualify the Dynare-mirror sentence: the equation mirror matches the
  fiscal structure, but the current quantitative calibration is the Julia
  \texttt{src/calibration.jl} block unless the Dynare parameter file is updated.
  \item In the policy-experiment paragraph, distinguish a wedge shock scale
  from a statutory-rate change.
\end{enumerate}

\end{document}
